% ============================================================================
% Fichier  : partie4/P04_C14_architecture.tex
% Titre    : Architecture de Sécurité
% Auteur   : MINKA MI NGUIDJOÏ Thierry Emmanuel
% Version  : 1.0 -- novembre 2026
% ============================================================================
% Base : brouillon « Livre IV » (principes, segmentation). Sources : Saltzer
% et Schroeder (1975) ; NIST SP 800-207 (architecture zéro confiance, 2020) ;
% CIS Controls v8.
% ============================================================================

\chapter{Architecture de Sécurité}
\label{chap:architecture}
\chaptermark{Architecture}

\epigraph{%
    « Supposer que l'attaquant est déjà à l'intérieur,\\
    et construire en conséquence. »%
}{Principe du cours}

\niveauChapitre{Intermédiaire}{Chapitres~\ref{chap:proprietes},
\ref{chap:menaces} et~\ref{chap:modeles}.}

\begin{boxobjectifs}
\begin{itemize}
    \item Appliquer la défense en profondeur à un système d'information.
    \item Découper un réseau en zones de confiance et définir les flux
          autorisés.
    \item Isoler l'administration et les sauvegardes.
    \item Expliquer les principes de l'architecture zéro confiance.
    \item Proposer une architecture cible pour un cas réel.
\end{itemize}
\end{boxobjectifs}

\bigskip

% ============================================================================
\section*{Question directrice : pourquoi un mot de passe volé a-t-il ouvert les quatre centres~?}
% ============================================================================

Le réseau de LiZbiZ était \emph{plat}~: une fois entré par l'accès
distant de Garoua, l'attaquant voyait les serveurs de Douala, de Yaoundé
et de Bafoussam, les sauvegardes et les contrôleurs de domaine. Rien ne
séparait le poste d'un technicien du cœur du système. Une architecture de
sécurité organise les séparations qui transforment une intrusion locale en
incident limité.

\medskip
\textbf{Plan du chapitre.}
\begin{enumerate}
    \item La défense en profondeur (section~\ref{sec:c14-profondeur}).
    \item Les zones et les flux (section~\ref{sec:c14-zones}).
    \item Isoler l'administration et les sauvegardes (section~\ref{sec:c14-admin}).
    \item L'architecture zéro confiance (section~\ref{sec:c14-zt}).
    \item Une architecture cible pour LiZbiZ (section~\ref{sec:c14-cible}).
\end{enumerate}

\separateur

% ============================================================================
\section{La Défense en Profondeur}
\label{sec:c14-profondeur}
% ============================================================================

\begin{boxdef}
\termecle{Défense en profondeur}~: superposition de mesures indépendantes,
de natures différentes, de sorte que la défaillance de l'une ne suffise
pas à compromettre l'ensemble.
\end{boxdef}

\begin{tableaucours}{L{3.0cm}L{9.2cm}}{%
    \tableauHeader{Couche} &
    \tableauHeader{Mesures}
}
Organisation & Politique, rôles, sensibilisation, contrats \\
Physique & Contrôle d'accès aux salles, alimentation, climatisation \\
Périmètre & Pare-feu, filtrage du courriel, accès distant authentifié \\
Réseau interne & Segmentation, filtrage entre zones, détection \\
Systèmes & Durcissement, correctifs, comptes distincts, EDR \\
Applications & Développement sécurisé, contrôle d'accès applicatif \\
Données & Chiffrement, classification, sauvegardes hors ligne \\
Détection et réponse & Journaux centralisés, surveillance, plan d'incident \\
\end{tableaucours}
\vspace{-0.8em}
\begin{center}{\small\itshape Tableau~14.1. Les couches de la défense en profondeur}\end{center}

L'Opération GAROUA a traversé toutes les couches parce que chacune
manquait ou était faible~: pas de second facteur au périmètre, pas de
segmentation, un correctif absent, des sauvegardes accessibles, aucune
détection.

% ============================================================================
\section{Les Zones et les Flux}
\label{sec:c14-zones}
% ============================================================================

\begin{boxdef}
\termecle{Zone de confiance}~: ensemble de systèmes de même sensibilité et
de même exposition, séparé des autres par un filtrage qui n'autorise que
les flux nécessaires. La \termecle{DMZ} (zone démilitarisée) accueille les
services exposés à Internet.
\end{boxdef}

\begin{tableaucours}{L{2.8cm}L{4.4cm}L{5.0cm}}{%
    \tableauHeader{Zone} &
    \tableauHeader{Contenu} &
    \tableauHeader{Flux autorisés (exemples)}
}
Internet & Tout & Vers la DMZ, en HTTPS seulement \\
DMZ & Portail de suivi, relais de messagerie, passerelle d'accès distant &
Vers les applications, sur des ports précis \\
Applications & Serveurs de l'application intégrée & Depuis la DMZ et les
postes, vers les bases \\
Données & Bases de données, partages de fichiers & Depuis les seules
applications \\
Postes & Postes des agences & Vers les applications, pas vers les bases \\
Administration & Bastion, outils d'administration & Vers toutes les
zones, en entrée depuis aucune \\
Sauvegarde & Serveurs et stockage de sauvegarde & En entrée, les seuls
flux de sauvegarde \\
\end{tableaucours}
\vspace{-0.8em}
\begin{center}{\small\itshape Tableau~14.2. Zones de confiance et flux d'une
entreprise comme LiZbiZ}\end{center}

\begin{boxalert}
\textbf{Violation typique.} Créer des réseaux virtuels séparés, mais
autoriser « tout » entre eux « pour que cela fonctionne ».

\smallskip
\textbf{Règle~:} une zone sans filtrage n'est qu'une étiquette. Chaque
flux entre zones est décrit (source, destination, port, motif,
propriétaire) dans une matrice, et tout le reste est refusé par défaut.
\end{boxalert}

% ============================================================================
\section{Isoler l'Administration et les Sauvegardes}
\label{sec:c14-admin}
% ============================================================================

Les comptes d'administration sont la cible de toute attaque avancée~:
qui les tient, tient le système. Trois règles les protègent.

\begin{enumerate}
    \item \textbf{Administrer depuis un poste dédié}, ou un \emph{bastion}
          qui concentre et enregistre les sessions, jamais depuis un poste
          qui navigue sur Internet ou lit des courriels.
    \item \textbf{Séparer les niveaux}~: les comptes qui administrent le
          domaine ne servent pas sur les postes de travail, où leurs
          identifiants pourraient être volés en mémoire.
    \item \textbf{Isoler les sauvegardes}~: une zone, des comptes et un
          stockage immuable qui ne dépendent pas du domaine de production
          (Chapitre~\ref{chap:modeles}, mesure M-13).
\end{enumerate}

% ============================================================================
\section{L'Architecture Zéro Confiance}
\label{sec:c14-zt}
% ============================================================================

L'architecture « périmétrique » fait confiance à ce qui est à l'intérieur
du réseau. L'architecture \termecle{zéro confiance}, décrite par le NIST
SP~800-207, n'accorde aucune confiance implicite à l'emplacement réseau~:
chaque accès à une ressource est décidé au moment de la demande, selon
l'identité, l'état du poste et le contexte, et pour la seule
ressource~\cite{nist800207}.

\begin{tableaucours}{L{4.6cm}L{7.6cm}}{%
    \tableauHeader{Principe} &
    \tableauHeader{Traduction}
}
Toute ressource est protégée & Même un serveur « interne » exige une
authentification \\
Toute communication est sécurisée & Chiffrement, quel que soit le réseau \\
Accès par session, au plus juste & Une session ne donne accès qu'à une
ressource \\
Décision dynamique & Identité, état du poste, heure, comportement \\
Surveillance continue & L'état des actifs et des accès est mesuré en
permanence \\
\end{tableaucours}
\vspace{-0.8em}
\begin{center}{\small\itshape Tableau~14.3. Les principes du zéro confiance
(d'après le NIST SP~800-207)}\end{center}

\begin{boxinfo}
\textbf{Une démarche, pas un produit.} Le zéro confiance se construit
progressivement~: authentification forte partout, inventaire des
ressources, segmentation fine, décisions d'accès centralisées. Une
entreprise qui n'a ni inventaire ni identités maîtrisées commence par là.
\end{boxinfo}

% ============================================================================
\section{Une Architecture Cible pour LiZbiZ}
\label{sec:c14-cible}
% ============================================================================

\begin{boxscenario}{Opération GAROUA. L'architecture après l'attaque}
\textbf{Proposition du RSSI.} (1)~Passerelle d'accès distant en DMZ, avec
double authentification et comptes nominatifs pour le prestataire, qui
n'ouvrent que les serveurs qu'il exploite~; (2)~séparation, dans chaque
centre, des zones applications, données, postes, administration et
sauvegarde~; (3)~filtrage entre centres limité à la réplication des
bases~; (4)~administration par un bastion qui enregistre les sessions~;
(5)~sauvegardes immuables dans une zone isolée, avec une copie hors
ligne~; (6)~journaux de toutes les zones centralisés et surveillés.

\medskip
\textbf{Analyse} (confiance moyenne). Avec cette architecture, le mot de
passe volé n'aurait ouvert que la passerelle (second facteur manquant),
et même avec lui, seulement les serveurs du prestataire, sans accès au
domaine ni aux sauvegardes.
\end{boxscenario}

\begin{boxmetier}{Administrateur systèmes et réseaux}
\textbf{Tenir la matrice des flux.} L'administrateur tient la matrice des
flux autorisés entre zones, la confronte chaque trimestre aux règles des
pare-feu, et supprime les règles sans propriétaire.
\end{boxmetier}

\begin{boxmetier}{RSSI}
\textbf{Faire valider l'architecture.} Le RSSI fait valider l'architecture
cible par la direction, avec un plan de migration par étapes et le risque
traité à chaque étape.
\end{boxmetier}

\separateur

\begin{boxresume}
\textbf{Ce qu'il faut retenir de ce chapitre~:}
\begin{itemize}
    \item \textbf{La défense en profondeur} superpose des mesures
          indépendantes, de l'organisation aux données.
    \item \textbf{Les zones de confiance} ne valent que par le filtrage des
          flux entre elles, refusé par défaut.
    \item \textbf{L'administration et les sauvegardes} s'isolent~: poste
          dédié ou bastion, niveaux séparés, stockage immuable.
    \item \textbf{Le zéro confiance} supprime la confiance implicite dans
          le réseau~; c'est une démarche progressive.
\end{itemize}
\end{boxresume}

\begin{boxinfo}
\textbf{Points de contrôle de l'élève.} Je sais~: décrire les couches de la
défense en profondeur~; découper un réseau en zones et rédiger la matrice
des flux~; expliquer pourquoi isoler l'administration~; énoncer les
principes du zéro confiance~; proposer une architecture cible.
\end{boxinfo}

\bigskip

\begin{boxexo}[title={\ding{45}\quad Exercice~14.1. Les couches \niveaubase}]
Pour chaque couche du Tableau~14.1, dites ce qui manquait chez LiZbiZ en
mars.
\end{boxexo}

\begin{boxexo}[title={\ding{45}\quad Exercice~14.2. Quelle zone~? \niveaubase}]
Placez dans une zone~: le serveur de messagerie relais~; la base des
clients~; le poste de la comptable~; le serveur de sauvegarde~; la console
de supervision.
\end{boxexo}

\begin{boxexo}[title={\ding{45}\quad Exercice~14.3. Réseau plat \niveaubase}]
Expliquez en cinq lignes pourquoi un réseau plat aggrave l'impact d'un
rançongiciel.
\end{boxexo}

\begin{boxexo}[title={\ding{45}\quad Exercice~14.4. La matrice des flux \niveaumoyen}]
Rédigez la matrice des flux (source, destination, port, motif,
propriétaire) entre les zones DMZ, applications et données de LiZbiZ.
\end{boxexo}

\begin{boxexo}[title={\ding{45}\quad Exercice~14.5. Le bastion \niveaumoyen}]
Quels avantages apporte un bastion d'administration~? Quels risques
crée-t-il, et comment les traiter~?
\end{boxexo}

\begin{boxexo}[title={\ding{45}\quad Exercice~14.6. Le prestataire \niveaumoyen}]
Concevez l'accès de SAHEL INFOGÉRANCE au système de LiZbiZ selon les
principes du zéro confiance.
\end{boxexo}

\begin{boxexo}[title={\ding{45}\quad Exercice~14.7. L'agence \niveaumoyen}]
Une agence de la BANQUE SANAGA a des guichets, un distributeur
automatique et une borne Wi-Fi pour les clients. Proposez ses zones et les
flux autorisés.
\end{boxexo}

\begin{boxexo}[title={\ding{45}\quad Exercice~14.8. Indépendance des mesures \niveaumoyen}]
Deux pare-feu du même fabricant, avec les mêmes règles, en série~:
est-ce de la défense en profondeur~? Justifiez.
\end{boxexo}

\begin{boxexo}[title={\ding{45}\quad Exercice~14.9. La migration \niveauexpert}]
Proposez un plan en quatre étapes pour passer de l'architecture actuelle
de LiZbiZ à l'architecture cible, avec le risque réduit à chaque étape.
\end{boxexo}

\begin{boxexo}[title={\ding{45}\quad Exercice~14.10. Rejouer l'attaque \niveauexpert}]
Rejouez l'Opération GAROUA étape par étape sur l'architecture cible, et
dites où elle s'arrête et qui la détecte.
\end{boxexo}

\bigskip

\begin{boxinfo}
\textbf{Pour aller plus loin.}
\begin{itemize}
    \item NIST SP~800-207, \emph{Zero Trust Architecture}~\cite{nist800207}.
    \item R.~Anderson, \emph{Security Engineering}~\cite{anderson2020}.
\end{itemize}
\end{boxinfo}

\begin{boxinfo}
\textbf{Transition.} L'architecture repose sur le réseau. Les
Chapitres~\ref{chap:reseaux1} et~\ref{chap:reseaux2} en étudient les
protocoles, les attaques et les équipements.
\end{boxinfo}
